%% =====================================================================
%%  T-01-01  The Five-Percent Premise
%%  -------------------------------------------------------------------
%%  One idea per file. Core is mandatory. Slots in canonical order.
%%  Max 700 words. No layout commands. No filler. New source material
%%  goes in sources/<BOOK>/T-01-01.tex -- never by rewriting this file.
%% =====================================================================
%% @kind: topic
%% @id: T-01-01
%% @title: The Five-Percent Premise
%% @subtitle: A minority operates; the majority is operated on
%% @chapter: c01
%% @order: 10
%% @status: stable
%% @sources: B1, B3
%% @updated: 2026-09-19

\topic{T-01-01}{The Five-Percent Premise}{A minority operates; the majority is operated on}

\begin{core}
Both sources open on the same claim: the ability to steer other people is not evenly distributed. A small minority exercises it deliberately and the large majority absorbs it. The minority is not smarter. It has simply decided that other people's behaviour is a variable it may act on, while the majority treats that behaviour as weather.
\end{core}
\begin{mechanism}
The asymmetry is maintained by an assumption rather than by a skill. Most people operate on a default rule --- that others are broadly like them, broadly honest, and broadly acting in good faith --- because that rule is cheap and usually correct enough. The practitioner suspends it. Once suspended, every interaction becomes readable in terms of interest, and reading interest is most of what the rest of this book teaches. The majority does not lose these exchanges because it is dull; it loses them because it is running a cooperative protocol against an opponent who is not.
\end{mechanism}
\begin{tells}
  \item You find out about an arrangement after it has been decided.
  \item Your reasons for refusing are answered one at a time, patiently, by someone who clearly knew them in advance.
  \item You concede things you did not intend to concede, and cannot locate the moment you did it.
  \item The other party's stated motive changes between tellings, while your own stays fixed.
\end{tells}
\begin{moves}
  \item Treat the distribution as a fact, not as a grievance.
  \item Assume any given counterpart may be in the minority until their behaviour proves otherwise.
  \item Convert the default rule from \emph{trust until shown otherwise} to \emph{verify while appearing to trust}.
\end{moves}
\begin{counters}
  \item Delay. Almost every technique in this book depends on momentum, and momentum dies overnight.
  \item Require writing. A person working you verbally will often decline to repeat it in a form you can keep.
  \item Ask the interest question out loud: \emph{what do you get if I say yes?} The answer, not the tone, is the information.
  \item Keep a record of your own concessions. Counting them removes the illusion that you chose each one freely.
\end{counters}
\begin{limits}
The premise is a posture, and postures have a cost. Held too tightly it produces a person who cannot accept help, cannot be surprised by generosity, and reads every kindness as a setup --- which is itself a form of being controlled, just by fear. Both sources acknowledge the trade. The usable version is situational: verify hard where money, commitment or reputation is at stake, and let the default rule stand where it does not.
\end{limits}

\sources{B1, B3}
\seesources
\seealso{T-01-02, T-04-01, T-06-04}
